\documentclass[11pt]{article}
\usepackage[margin=0.8in]{geometry}
\usepackage{amsmath,microtype,hyperref}
\newcommand{\guideheading}[1]{\par\medskip\noindent{\large\bfseries #1}\par\smallskip}
\begin{document}
\begin{center}
{\Large Guide: Checking What a Residue Table Really Predicts}\par\medskip
Collatz proof project\quad---\quad September 5, 2026
\end{center}
\guideheading{Why this matters for the proof}
A table that predicts the next residue could help organize a global Collatz
argument. But its input must retain enough information to determine its
output. Division by two can make the next output depend on a bit that the
input table discarded. We checked this issue before adopting a proposed
one-bit mixing route from the literature.

\guideheading{A concrete failure}
The numbers 3 and 35 have the same remainder modulo 32: both leave 3.
Both are eligible for the odd shortcut step, which multiplies by three,
adds one, and divides by two. The results are
\[
 3\longmapsto5,\qquad 35\longmapsto53.
\]
Their output remainders modulo 32 are different: 5 and 21. A single
transition attached to input remainder 3 cannot describe both.
Lean proves that no function of this input remainder alone works for all
such gap steps.

\guideheading{The missing information is identifiable}
For every natural $q$, Lean checks
\[
 T(3+32q)=5+48q.
\]
The parity of $q$ selects the output remainder. Recording the input modulo
64 would retain that bit. This is a precision requirement, not evidence
that modular methods in general cannot work.

\guideheading{Compressing several divisions needs care too}
A compressed odd step divides out every factor of two after applying
$3n+1$. Inputs 29 and 61 agree modulo 32, yet their compressed outputs are
11 and 23. Those outputs disagree even modulo eight. A table using the
smaller representative 29 cannot automatically stand in for 61.

\guideheading{What the count check says}
Among the canonical burst inputs below 32, two compressed images are
3 modulo eight and one is 7 modulo eight. Lean verifies those counts.
Their signed difference is positive one. This disagrees with the signed
formula at depth five in the reviewed source, while agreeing with an
absolute difference of one. The technical note identifies the exact
version and statements being checked.

\guideheading{What we can conclude}
These checks prevent a particular loss of information from entering the
formal proof. They do not evaluate every claim in the source, exclude a
corrected argument, or establish Collatz. A repaired residue description
would still need a proof about the behavior of every relevant orbit.

\guideheading{Verification files}
Build \texttt{Collatz.ResidueMapAudit} from the repository's \texttt{lean/}
directory. The accompanying technical note provides the exact equations,
source reference, and limits of the conclusion. The project axiom audit
includes all five declarations.
\end{document}
